\subsection{Operations without local monodromy of the Poincaré groupoid}

Lemma 1. --- Let B be a locally arcwise connected topological space, let \(p \colon E \to B\) and \(p' \colon E' \to B\) be étale and separated maps satisfying the path lifting property (III, p. 302). Let \(g \colon E \to E'\) be a mapping such that \(p' \circ g = p\) . We assume that \(g\) is compatible with the canonical operations of the groupoid \(\varpi(B)\) on \(E\) and \(E'\) . Then, \(g\) is continuous.

Let c be a path in E; set \(c' = g \circ c\) and let us prove that the map \(c'\) is continuous. Let d denote the path \(p \circ c\) in B and, for every \(t \in I\) , denote by \(d_t\) the path \(s \mapsto d(st)\) . For all \(t \in I\) , we have \(c(t) = c(0) \cdot d_t\) (III, p. 304, remark 3), whence \(c'(t) = c'(0) \cdot d_t\) by the hypothesis made on g, which proves that \(c'\) is a path lifting the path d, by this same remark. This proves that the map g is continuous by arcs. Since the space E is locally arc-connected (III, p. 261, cor. 2), the map g is continuous (III, p. 269, corollary of prop. 13).

Let B be a topological space. Consider an operation \(\varphi = (\varphi_{a,b})_{(a,b)\in \mathrm{B}\times \mathrm{B}}\) of the groupoid \(\varpi (\mathrm{B})\) on a set E, with respect to a mapping \(p\colon \mathrm{E}\to \mathrm{B}\) . One says that \(\varpi (\mathrm{B})\) acts without monodromy on E (cf. II, p. 168) if for every \(b\in \mathrm{B}\) and every loop class \(\gamma \in \pi_1(\mathrm{B},b)\) the action of \(\gamma\) on the fiber \(\mathrm{E}_b\) is trivial. If B is path-connected, it suffices that this be the case for one point of B (loc. cit.). We shall say that the operation \(\varphi\) of the groupoid \(\varpi (\mathrm{B})\) is without local monodromy if every point of B has a neighborhood V such that \(\varpi (\mathrm{V})\) acts without monodromy on the set \(\mathrm{E_V} = p^{-1}(\mathrm{V})\) with respect to the map \(p_{\mathrm{V}} = p\mid p^{-1}(\mathrm{V})\) .

Remark. --- Let B be a locally path-connected topological space, and suppose that every point b of B has a neighborhood V such that the image of \(\pi_{1}(V,b)\) in \(\pi_{1}(B,b)\) is reduced to the identity element (cf. IV, p. 340, def. 2). Then, every operation of the groupoid \(\varpi(B)\) is without local monodromy.

Indeed, consider a set E, a mapping \(p \colon \mathrm{E} \to \mathrm{B}\) and an operation \(\varphi\) of the groupoid \(\varpi(\mathrm{B})\) on E with respect to \(p\) . Let \(a\) be a point of B, and let V be a neighborhood of \(a\) such that the image of \(\pi_1(\mathrm{V}, a)\) in \(\pi_1(\mathrm{B}, a)\) is reduced to the identity element. Let U be a path-connected neighborhood of \(a\) contained in V. Let \(b\) be a point of U and let \(\gamma \in \pi_1(\mathrm{U}, b)\) . Let \(\delta\) be the class of a path joining \(a\) to \(b\) in U. Then, \(\delta\gamma\delta^{-1}\) is the class of a loop at \(a\) in U; its image in \(\pi_1(\mathrm{B}, a)\) is therefore trivial. We thus have \(\varphi_{a,a}(\delta\gamma\delta^{-1}) = \mathrm{Id}_{\mathrm{E}_a}\) , whence \(\varphi_{b,b}(\gamma) = \mathrm{Id}_{\mathrm{E}_b}\) . This proves that the operation of \(\varpi(\mathrm{U})\) on the U-space \(\mathrm{E_U}\) is without monodromy. Consequently, the operation \(\varphi\) is without local monodromy.

PROPOSITION 3. --- Let B be a locally path-connected topological space and let E be a set equipped with an operation \(\varphi\) of the groupoid \(\varpi(B)\), without local monodromy, relative to a mapping \(p: E \to B\). There then exists on E a unique topology for which the following conditions are satisfied:

\begin{enumerate}
\def\labelenumi{(\roman{enumi})}
\item
  The set E equipped with this topology and with the map p is a covering of B;
\item
  The canonical action of \(\varpi(\mathrm{B})\) on this covering is identical to the action \(\varphi\) .
\end{enumerate}

The uniqueness of such a topology follows from lemma 1 of III, p. 312, where one takes for g the identity map of E. To prove its existence, one may moreover assume that B is connected and nonempty.

Suppose first that the operation \(\varphi = (\varphi_{a,b})_{(a,b)\in \mathrm{B}\times \mathrm{B}}\) of the groupoid \(\varpi (\mathrm{B})\) on the set E, relative to \(p\) , is without monodromy.

Let \(a\) be a point of B. For every point \(b \in \mathrm{B}\) , there exists a path \(c\) in B joining \(a\) to \(b\) . If \(\gamma \in \varpi_{a,b}(\mathrm{B})\) denotes the class of the path \(c\) , the bijection \(\varphi_{a,b}(\gamma) \colon \mathrm{E}_a \to \mathrm{E}_b\) is independent of the path \(c\) , since the operation is without monodromy; we denote by \(f_{a,b}\) this bijection. The map \(\Phi_a \colon \mathrm{B} \times \mathrm{E}_a \to \mathrm{E}\) defined by \((b,x) \mapsto f_{a,b}(x)\) is a bijection; the inverse bijection maps \(x \in \mathrm{E}\) to the pair \((p(x), f_{p(x),a}(x))\) . Endow the set \(\mathrm{E}_a\) with the discrete topology, so that the B-space \(\mathrm{B} \times \mathrm{E}_a\) is a trivial covering of B. By transport of structure, the bijection \(\Phi_a\) endows E with a topology making it a covering of B.

Let us prove that the canonical operation of \(\varpi(\mathrm{B})\) on E is identical to the operation \(\varphi\) . Let \(x\) be a point of \(\mathrm{E}_a\) and \(b\) a point of B. Let \(c\) be a path in B joining the point \(a\) to the point \(b\) ; the map \(t \mapsto \Phi_a(c(t), x)\) is then a path in E joining the point \(x\) to the point \(\Phi_a(b, x) = f_{a,b}(x)\) . If \(\gamma \in \varpi_{a,b}(\mathrm{B})\) denotes the class of the path \(c\) , we thus have \(x \cdot \gamma = f_{a,b}(x)\) . This shows that the canonical action of \(\varpi(\mathrm{B})\) on the covering E and the operation \(\varphi\) coincide on the classes of paths of origin \(a\) . Since these classes generate the groupoid \(\varpi(\mathrm{B})\) , the two operations are equal.

Now let us treat the general case. Let \(\mathcal{B}\) be the set of open subsets V of E such that \(\varpi(\mathrm{V})\) acts without monodromy on \(p^{-1}(\mathrm{V})\) . By hypothesis, the elements of \(\mathcal{B}\) cover B. By the preceding, there exists for every \(\mathrm{V} \in \mathcal{B}\) a unique topology on the set \(p^{-1}(\mathrm{V})\) such that \((p^{-1}(\mathrm{V}), p_{\mathrm{V}})\) makes it a covering of V and such that the canonical action of \(\varpi(\mathrm{V})\) on this covering coincides with the action induced by \(\varphi\) .

Let V and \(V' \in \mathcal{B}\) . The topology of \(p^{-1}(V \cap V')\) induced by the topology of \(p^{-1}(V)\) (resp. of \(p^{-1}(V')\) ) defined above indeed makes a covering of \(V \cap V'\) on which the canonical operation of \(\varpi(V \cap V')\) is induced by the operation \(\varphi\) . These topologies therefore coincide with that of \(p^{-1}(V \cap V')\) . There then exists a unique topology on E inducing on each \(p^{-1}(V)\) the topology previously defined (cf. I, §2, p. 16).

When E is equipped with this topology, the map p is continuous and the B-space E is a covering. The canonical action of \(\varpi(\mathrm{B})\) on this covering coincides with the operation \(\varphi\) on the classes of paths whose image is contained in one of the open sets of \(\mathcal{B}\) . By Lemma 4 of

III, p. 272, these classes generate the groupoid \(\varpi(B)\) . It follows that these two operations are equal (II, p. 167).

COROLLARY. --- Let B be a connected and locally arcwise connected topological space, let (E,p) be a B-space whose projection p is étale, separated, and has the path lifting property. If the canonical operation of \(\varpi(\mathrm{B})\) on the set E, relative to p, is without local monodromy, then E is a covering of B.

Endow E with the canonical action of \(\varpi(B)\) defined by the lifting of paths (III, p. 303, no. 3). There exists a topology on E for which conditions (i) and (ii) of Prop. 3 are satisfied. This topology coincides with the given topology on E by Lemma 1 of III, p. 312.

